% !TeX root = ../papers/nc-p3.tex
\section{対角交換子零の疎な四元群窓と二つの正規平方}
\label{sec:marked-klein-zero-bracket-boundary}
\begin{theorem}\label{thm:marked-klein-zero-bracket-boundary}
定理\ref{thm:marked-klein-mixed-supports}の四元群次数で共通差を $\delta=0$ とし、
少なくとも一成分が零となる場合を考える。標準14641型quadratic Koszul左右AS正則になる台は、
軸の共通置換と四位置の巡回を除き両側とも $(0,15,15)$ に限る。
この台は $x,b\in k^*$ による下記表示で覆われ、その全点が正則である。
特に $x=b$ で二組の平方を結ぶ行列が階数1に落ちても例外は生じない。
\end{theorem}

\begin{proof}[四位置を通る単項式商による一括排除]
各局所台は $2\times2$ 行列の階数が高々1なので空または二つの部分集合の直積であり、十通りある。
前節の十二の非対角式で積の零・非零を一致させると47492対、
共通軸置換と巡回で2038軌道である。
これは台の必要条件による列挙なので、非実現な台を先に除くことは不要である。

全位置で星文字を零にして、残りの関係に現れる二文字も個別に零とした単項式商を取る。
各位置で $jj$ は残る。相異なる $j,k$ と残りの $i$ に対し、$jk$ がその位置で残る必要十分条件は
$(G_j)_{ki}=(Q_k)_{ji}=0$ である。実際、後ろ二位置の $(0,i),(i,0)$ による六縮約が全関係の基底となる。
三頂点有向グラフに、一周期内で少なくとも一度残る全ての $j\to k$（$j\ne k$）を加える。
閉路があれば、各文字で留まる遷移を用いて、閉路を一周して開始位置と開始文字に戻れる。
同じ長さの一定語とこの周回を任意に連接できるため、経路数は指数増大する。
二つの三つ組について十二個の積の真理値を突き合わせ、各軌道の上のグラフの推移閉包を取ると、
閉路を持たないのは両側が $(0,15,15)$ の場合と両側が $(15,15,15)$ の場合だけである。
後者は疎ではない。\path{computations/verify_klein_zero_bracket_boundary.py} が全台を再生成してこの判定を行う。
\end{proof}

\begin{proof}[残る台の二変数表示]
軸を $1,2,3$ とし $Q_1=P_1=G_1=F_1=0$ である。第2,3軸の $Q_i^\pm$ は、
第一枠の二つの独立な隣接対角比の平方、および第二枠の二倍率で全て1にできる。
$P_i$ はそれぞれ $Q_i$ の非零倍である。共通倍率で $P_2=Q_2$, $P_3=xQ_3$ にする。
$G_2,G_3$ の正成分は第一枠の共通倍率および未使用の第二枠第1軸の倍率でともに1にできる。
残る非対角四式は
\[
 (G_2^+,G_2^-)=(1,b),\quad (G_3^+,G_3^-)=(1,x/b),\qquad F_i=-G_i
\]
を与える。実際 $F_i=k_iG_i$ と書けば正成分同士の二式が $k_2=k_3=-1$、
負成分同士の式が $G_2^-G_3^-=x$ である。他の式は恒等的に成り立つ。
\end{proof}

\begin{proof}[正規列と全次数の自由基底]
文字を $0,1,2,3$ で記し、奇数位置でのみ $0,1$ および $2,3$ を交換する。
以後 $ii$ はその位置からの二文字を表す。$i=0,1$ と $j=2,3$ の間の四関係は
$ji=c_{p,ij}ij$、他の二つは $ii=D_{p,i2}\,22+D_{p,i3}\,33$ ($i=0,1$) となる。
$D_{p,2*},D_{p,3*}$ は単位行とする。全 $c$ と $D_{p,0j},D_{p,1j}$ は非零Laurent単項式である。
例えば位置0の二つの行は $(-1,-x/b),(-1,-1)$ である。各位置の係数は上の四組から直接得られる。

$z_2=22,z_3=33$ を次数2の追加文字とする。これらは各一文字を非零スカラー倍で通過する。
実際 $i=0,1$ では二回の交差交換で $i z_j=(c_{p,ij}c_{p+1,ij})^{-1}z_j i$ である。
$i=2,3$ は $i z_i=z_i i$ と、$i$ を $00$ の左から右へ動かす式を使う。
左の平方係数は $D_{p+1,0*}$、右は $c_{p,0i}c_{p+1,0i}D_{p,0*}$。
第 $i$ 単位行と $D_{p+1,0*}$ の行列式は非零単項式なので、二行だけで
通過行列が決まり、検算すると対角で可逆である。これが $x=b$ でも有効な理由である。
$z_3 z_2$ も二回の通過で $z_2 z_3$ の非零倍となる。

書換えは四平方を $z$ の一次結合へ、混合対を $ij$ ($i<2\le j$) へ、$z$ を先頭へ動かして
$z_2,z_3$ の順とする。全て元の関係の帰結であり、逆に $z$ を消去すれば元の表示になる。
平方では総文字数が減り、他の書換えでは $z_2,z_3$, 第一組、第二組という順に対する転倒数が減る。
二文字書換えの重なりは各位置32個で、記載のコードで左右の正規化を Laurent 多項式として照合すると
全て一致する。従って基底は $z_2^a z_3^b$、$0,1$ の交代語、$2,3$ の交代語を順に並べた全語である。
各組の交代語の母関数は $(1+t)/(1-t)$、二つの $z$ は次数2なので
各始点のHilbert級数は $(1-t)^{-4}$ である。両側で $z_2$ およびその商での $z_3$ は正規非零因子である。

最後に定理\ref{thm:klein-rational-period-four}の線型複体の議論を適用する。
折り畳んだ環で全位置の同じ $z_j$ を足せば、通過式と自由基底から正規正則列となる。
従ってどの頂点単純から環への $\operatorname{Hom},\operatorname{Ext}^1$ も両側で消える。
実際最初の正規元による商の二項分解での変換環系列では、一段目のExtは商環へのHomに
正規化自己同型と次数移動を施したものになり、二番目の非零因子がこれを消す。
四テンソルの縮約による $1,4,6,4,1$ 複体は両端で互いの双対であり、
隣接切断の階数6、星成分による一文字縮約の単射性と二次表示から始端の二写像は表示を与える。
上の二消滅で逆端も完全になり、Hilbert級数のEuler等式から中央も完全となる。
左右とも線型な最小分解とその双対の最終単純を得るので、標準Koszul左右AS正則である。
\end{proof}
